%%=============================================================================
%% Voorwoord
%%=============================================================================

\chapter*{\IfLanguageName{dutch}{Woord vooraf}{Preface}}%
\label{ch:voorwoord}

%% TODO:
%% Het voorwoord is het enige deel van de bachelorproef waar je vanuit je
%% eigen standpunt (``ik-vorm'') mag schrijven. Je kan hier bv. motiveren
%% waarom jij het onderwerp wil bespreken.
%% Vergeet ook niet te bedanken wie je geholpen/gesteund/... heeft

Dit onderzoek heeft me doen ingezien waar software ontwikkeling naartoe aan het gaan is. Deze AI ondersteunde workflow is inzetbaar op veel vlakken van ontwikkeling niet enkel op testing. Door deze bachelorproef te voeren heb ik de kans gekregen om echt te duiken in de nieuwe tools die de industrie aan het omarmen is.
\bigskip
\\
In het begin stond ik twijfelachtig tegenover het gebruik van AI tools, doordat ik er van uit ging dat basisprincipes herhalen en moeilijkere problemen op lossen de weg waren om mij verder te ontplooien. Begrijp mij niet verkeerd, als je niet kan coderen gaat het moeilijk worden. Maar om een doeltreffende ontwikkelaar te zijn moet ik niet alleen omarmen wat er nu is, maar ook de rol van een ontwikkelaar verbreden.
\bigskip
\\
Namelijk de rol is aan het veranderen. Het gaat niet enkel meer over goede code te kunnen schrijven, maar ook over het kunnen documenteren en het gieten van taken in processen die geautomatiseerd kunnen worden. Hiervoor heb ik sterk geleund op de analyse vakken van HOGENT. Eenmaal de vereisten zijn opgesteld, de richtlijnen gedefinieerd zijn, implementeert de agent het werk. En om kwaliteitsvolle output te hebben moet de uitvoerder zijn basisprincipes beheersen.
\bigskip
\\
Deze ervaring maakt me zeer tevreden dat ik geleid en gesteund was om dit onderzoek uit te voeren.
\bigskip
\\
Graag wil ik een aantal mensen bedanken die mij geholpen hebben om dit op het juiste pad te krijgen. Namelijk J. Straetemans die me heeft geholpen om meer product minded te denken en kwalitatief onderzoek te plegen. Verder mijn co-promoter M. Soufi voor de ondersteuning, en natuurlijk ook mijn promotor M. Van Audenrode, die mij geholpen heeft kritische vragen te stellen over het onderwerp, het onderzoek in de juiste vorm te gieten en haar waardevolle feedback om dit te vervolledigen.
\bigskip
\\
Ook wil ik mijn familie en vrienden bedanken om mij de nodige ruimte en ondersteuning te geven. 